\section{API Specification}

The communication layer is built on gRPC/Protobuf. The core service is \texttt{DemocritusScheduler}.

\subsection{Service Definition}
The \texttt{DemocritusScheduler} service facilitates the coordination between the Master (Scheduler) and the Workers. It exposes three main RPC methods:
\begin{itemize}
    \item \texttt{RegisterWorker}: Handshake method for workers to join the grid.
    \item \texttt{GetTask}: Workers pull work from the scheduler.
    \item \texttt{SubmitResult}: Workers push the computation results back.
\end{itemize}

\subsection{TaskPayload Structure}
The \texttt{TaskPayload} message is critical for scientific reproducibility. It contains:
\begin{itemize}
    \item \texttt{task\_id} (string): Unique identifier for the job slice.
    \item \texttt{seed} (int64): The seed for the Pseudo-Random Number Generator (PRNG). By transmitting this from the master, we ensure that re-running Task X always yields result Y, regardless of which worker computes it.
    \item \texttt{iteration\_count} (int64): The number of simulation steps (e.g., random walk steps).
    \item \texttt{simulation\_type} (string): Identifies the physics model, e.g., ``random\_walk\_3d''.
\end{itemize}
